\documentclass[12pt, oneside, openany]{book}

\usepackage{mathptmx} % Contiene una fuente similar a Times New Roman

\usepackage[spanish, es-tabla]{babel} % Permite escritura en castellano
\usepackage[utf8]{inputenc} % Permite utilizar caracteres UTF8

\usepackage{graphicx} % Para la inclusión de gráficos e imágenes
\graphicspath{ {images/} } % Ruta para buscar las imágenes
\usepackage[a4paper,top=30mm,left=30mm,right=25mm,bottom=25mm,headheight=20mm]{geometry} % Configuración de los margenes de la página

% Paquetes para que funcione el formato.
\usepackage{titlesec}
\usepackage{setspace}
\usepackage{ragged2e}
\usepackage{fancyhdr}
\usepackage{lastpage}
\usepackage{stackengine}
\usepackage{array}
\usepackage[hidelinks]{hyperref}
\usepackage{enumitem}
\usepackage{float}
\usepackage{hypcap}
\usepackage{caption}
\usepackage{fancyvrb}

\usepackage{hyperref} % Paquete para que las referencias funcionen, y permite introducir links
\usepackage{xcolor} % Paquete para trabajar con colores (fondo de celdas, color del texto...)
% Se define un color gris desde su código RGB
\definecolor{gris}{RGB}{220,220,220}

\setcounter{secnumdepth}{3} % Para permitir numerar las sub-subsecciones

% Modifica el nombre de los índices al castellano
\addto\captionsspanish{
  \renewcommand{\contentsname}{Índice de contenido}
  \renewcommand{\listfigurename}{Índice de figuras}
  \renewcommand{\listtablename}{Índice de tablas}
  \renewcommand{\lstlistingname}{Listado}
  \renewcommand{\lstlistlistingname}{Índice de listados}
}

% Formateo de los nombres de los apartados:
\titleformat{\chapter}[block]
{\normalfont\Huge\bfseries\singlespacing}{\thechapter.}{1em}{\Huge}
\titlespacing*{\chapter}{0pt}{-62pt}{0pt}

\titleformat{\section}[block]
{\normalfont\Large\bfseries}{\thesection.}{4pt}{\Large}
\titlespacing*{\section}{0pt}{\baselineskip}{0pt}

\titleformat{\subsection}[block]
{\normalfont\large\bfseries}{\thesubsection.}{4pt}{\normalsize\large}
\titlespacing*{\subsection}{0pt}{0pt}{0pt}

\titleformat{\subsubsection}[block]
{\normalfont\normalsize\bfseries}{\thesubsubsection.}{4pt}{\normalsize}
\titlespacing*{\subsubsection}{0pt}{0pt}{0pt}

\def\tablename{Tabla}

%% Variables para portada y cabeceras
%% Cambiar los valores para cada documento!!!
\def\title{Explotación, integración y visualización de múltiples fuentes de datos mediante un Data Lake}
\def\shortTitle{ETL de datos heterogéneos mediante un Data Lake}
\def\subject{Lenguajes y Sistemas Informáticos}
\def\author{Mier Montoto, Juan Francisco}
\def\date{julio 2024}
\def\org{Escuela Politécnica de Ingeniería de Gijón}
\def\area{Grado en Ingeniería Informática en Tecnologías de la Información}
\def\tutorOne{D. Augusto Alonso, Cristian}
\def\tutorTwo{D. Morán Barbón, Jesús}
\def\tutorThree{D. Vázquez Faes, Eduardo}
\def\tfgId{???}

\def\ORG{\expandafter\MakeUppercase\expandafter{\org}}
\def\AREA{\expandafter\MakeUppercase\expandafter{\area}}
\def\SUBJECT{\expandafter\MakeUppercase\expandafter{\subject}}


\captionsetup{justification=centering}
\setlength{\headheight}{65pt}

\fancyhf{}
\fancyhead[L]{\includegraphics[height=16mm]{style/square.png}
  \hspace{1em} \Longstack[l] {
    \textbf{\SUBJECT} \newline
    \textbf{\shortTitle}}
  \newline \leftmark{}
}
\fancyhead[R]{\bfseries{Hoja \hyperlink{toc}{\thepage}~de~\pageref{LastPage}}}
\fancyfoot[C]{\author}
\renewcommand{\headrulewidth}{0pt} % default is 0pt
\renewcommand{\footrulewidth}{0.4pt} % default is 0

\fancypagestyle{plain}{%
  \fancyhf{}
  \fancyhead[L]{\includegraphics[height=16mm]{style/square.png}
    \hspace{1em} \Longstack[l]{
      \textbf{\SUBJECT} \newline
      \textbf{\shortTitle}}}
  \fancyhead[R]{\bfseries{Hoja \hyperlink{toc}{\thepage}~de~\pageref{LastPage}}}
  \fancyfoot[C]{\author}
  \renewcommand{\headrulewidth}{0pt} % default is 0pt
  \renewcommand{\footrulewidth}{0.4pt} % default is 0pt
}

\newcommand*{\fullref}[1]{\textit{\hyperref[{#1}]{\ref*{#1} \nameref*{#1}}}}

\pagestyle{fancy}

\restylefloat{table}
\definecolor{backcolour}{rgb}{0.95,0.95,0.92}
\definecolor{darkgreen}{rgb}{0.0, 0.5, 0.0}

\lstdefinestyle{default}{
  basicstyle=\ttfamily\footnotesize,
  breakatwhitespace=false,
  breaklines=true,
  captionpos=b,
  keepspaces=true,
  numbers=left,
  numbersep=5pt,
  numberstyle=\tiny\color{black},
  showspaces=false,
  showstringspaces=false,
  showtabs=false,
  tabsize=2,
  backgroundcolor=\color{backcolour},
  postbreak=\mbox{\textcolor{red}{$\hookrightarrow$}\space},
}

\lstset{style=default}

\lstdefinestyle{yaml}{
  basicstyle=\ttfamily\footnotesize\color{darkgreen}\bfseries,
  rulecolor=\color{black},
  string=[s]{'}{'},
  stringstyle=\color{darkgreen},
  comment=[l]{:},
  commentstyle=\color{blue},
  morecomment=[l]{-},
  numbers=left,
  numbersep=5pt,
  numberstyle=\tiny\color{black},
  backgroundcolor=\color{backcolour},
  breaklines=true,
  captionpos=b,
  keepspaces=true,
  showspaces=false,
  showstringspaces=false,
  showtabs=false,
  tabsize=2,
  breakatwhitespace=true,
  postbreak=\mbox{\textcolor{red}{$\hookrightarrow$}\space},
}



\begin{document}

\rmfamily % Fuente tipo Romana

% Portada de la memoria
\begin{titlepage}
    \centering
    \bfseries {
        \null{}
        \vspace{0cm}
        \begin{table}[h]
            \centering
            \begin{tabular}{m{10cm} m{1cm} m{3cm}}
                \vspace{0.2cm}
                \includegraphics[width=86mm]{style/full.png} &  & \vspace{1.52mm} \includegraphics[width=23mm]{style/square.png} \\
            \end{tabular}
        \end{table}

        \vspace{3\baselineskip}

        \Large{\ORG{} \\ \vspace{3\baselineskip}}
        \large {
            \AREA{} \\ \vspace{3\baselineskip}
            \subject{} \\ \vspace{2\baselineskip}

            TRABAJO FIN DE GRADO/MÁSTER Nº \tfgId{} \vspace{\baselineskip} \\
            \title{} \\ \vspace{1\baselineskip}

            \author{} \\ \vspace{1\baselineskip}
            TUTORES:\\
            \tutorOne{} \\
            \tutorTwo{} \\
            \tutorThree{} \\ \vspace{\baselineskip}

            \vspace{2\baselineskip}
            FECHA:\@ \date{}
        }
    }
\end{titlepage}


% Índice de contenido
\addcontentsline{toc}{chapter}{Índice de contenido} % Añade la referencia al índice de contenido
\hypertarget{toc}{}
\tableofcontents
\newpage

% Índice de figuras
\addcontentsline{toc}{chapter}{Índice de figuras}  % Añade la referencia al índice de contenido
\hypertarget{lof}{}
\listoffigures

\justify{} % Texto justificado
\setlength{\parskip}{\baselineskip} % Separación entre párrafos de 1 linea
\onehalfspacing{}

%% El contenido de la memoria, dividido en capítulos:
\chapter{Introducción}\label{chap:intro}
En este capítulo se presenta una introducción al trabajo de desarrollo de software realizado,
proporcionando un contexto general y estableciendo el escenario para los capítulos siguientes.
Se discutirán los antecedentes y la motivación detrás de este trabajo, la finalidad del proyecto,
y se proporcionará una breve descripción de la empresa en la que se ha desarrollado este trabajo.
Este capítulo tiene como objetivo proporcionar una visión general del proyecto y establecer las
bases para los capítulos detallados que siguen.

\section{Antecedentes}\label{sec:antecedentes}
Hoy en día, nos encontramos en una era donde la generación y almacenamiento de datos crece
exponencialmente\footnote{\url{https://www.statista.com/statistics/871513/worldwide-data-created/}},
un hecho que se ve reflejado en el ámbito empresarial. La diversidad de fuentes y formatos de
estos datos introduce una complejidad significativa en su manejo, conocida como \textit{heterogeneidad}
\footnote{\url{https://www.sciencedirect.com/topics/computer-science/data-heterogeneity}}, siendo las
bases de datos, archivos de registros y APIs las fuentes más habituales.

El término \textit{big data} describe este fenómeno de acumulación masiva de datos, cuya magnitud y
complejidad sobrepasan las capacidades de los métodos de procesamiento convencionales. El \textit{big data}
por tres características principales: volumen, variedad y velocidad - su adecuada gestión y análisis
pueden otorgar ventajas competitivas significativas a las empresas, tales como el descubrimiento de
patrones ocultos, identificación de nuevas oportunidades de mercado y optimización de procesos de toma
de decisiones.

Uno de los procesos que permite la extracción de esta información es la pirámide DIKW, \cite{enwiki:1211227190}
es un modelo que describe la relación entre los datos, la información, el conocimiento y la sabiduría.
Según este modelo, los datos son la materia prima de la información, que a su vez es la materia prima
del conocimiento, que a su vez es la materia prima de la sabiduría. Una organización sin los procesos
adecuados para la gestión y análisis de estos datos, se enfrenta a importantes desafíos, como la dificultad
para identificar patrones y tendencias, la toma de decisiones incorrectas y la pérdida de oportunidades de
negocio. Por otro lado, una organización que logre extraer información valiosa de sus datos, podrá mejorar su
eficiencia, aumentar su competitividad y adaptarse mejor a un entorno empresarial en constante cambio.

La evolución tecnológica ha propiciado el desarrollo de innovadoras herramientas y metodologías
diseñadas para enfrentar estos desafíos. Entre ellas, los \textit{data lakes} (o \emph{lagos de
información}) se destacan por su capacidad para consolidar vastos volúmenes de datos heterogéneos,
facilitando su posterior análisis y aprovechamiento de manera más efectiva.

Sin embargo, a pesar de que existen herramientas de almacenamiento, el proceso de integración,
visualización y análisis de estos datos es una tarea desafiante, ya que requiere de una gran
cantidad de recursos y de un tiempo de desarrollo considerable del que, normalmente, no se dispone
en el ámbito empresarial.

Con la ingesta masiva de datos, se presentan nuevos problemas a la hora de analizar y obtener
información de ellos:

\begin{itemize}
	\item Sin la necesaria automatización y correcta aplicación de los procesos ETL,
		al tratarse de un crecimiento exponencial de los datos y, por lo tanto, de la fuerza
		de trabajo necesaria para manejarla, los resultados del análisis pueden ser incorrectos,
		lo que deriva en errores y decisiones de negocio equivocadas que impactan negativamente
		en la empresa.
	\item La heterogeneidad de los datos, tanto en formato como en origen, dificulta su
		consolidación y análisis.
	\item La masificación de información impide el análisis manual de los mismos, requiriendo resúmenes
		estadísticos o representaciones gráficas como \textit{dashboards} para su correcta interpretación.
		La visualización de datos es una técnica que permite representar la información de manera visual,
		para facilitar su análisis y comprensión. La visualización de datos es una parte importante del
		proceso de análisis de datos, ya que permite identificar patrones, tendencias y anomalías en los
		datos de forma más rápida y sencilla.
\end{itemize}

\newpage{}
\section{Motivación}\label{sec:motivacion}
Okticket, como el resto de empresas, se enfrenta a la necesidad de gestionar y analizar grandes
volúmenes de datos, provenientes de múltiples fuentes y en diferentes formatos. La correcta gestión
y análisis de estos datos es fundamental para la toma de decisiones y para la mejora de los procesos
internos de la empresa.

En la actualidad, la empresa dispone de una gran cantidad de datos que se encuentran en
diferentes formatos y en diferentes ubicaciones, lo que dificulta su análisis y explotación.
Por otra parte, se depende de la consulta manual o de servicios de terceros para poder analizar
estos datos, lo que supone un coste adicional.

El proyecto surge de la necesidad de la empresa de extraer información y conocimiento de las
múltiples y heterogéneas fuentes de datos de las que se disponen, tanto internas (e.g.~bases de
datos, archivos de registros, APIs, entre otros), como externas (e.g. APIs o datos de webs de
terceros, datos de fuentes públicas\ldots).

Además del uso interno, la empresa también quiere ofrecer a sus clientes la posibilidad de
consultar estos datos de forma visual y sencilla, para que puedan analizarlos y explotarlos de
forma autónoma, lo que supondría un valor añadido para los mismos.

\newpage{}
\section{La empresa}\label{sec:empresa}
Okticket es una startup nacida en Gijón en 2017 cuyo producto principal es un servicio software que escanea
automáticamente de tickets y notas de gastos lo que permite reducir los costes y el tiempo que invierten
las empresas en contabilizar y manejar los gastos de viaje de los profesionales.

La empresa tienen su suede principal en el Parque Tecnológico de Gijón, aunque cuenta con un número
de sedes creciente en varios países, como Francia, Portugal o, más recientemente, México. En esta
oficina principal se encuentran los departamentos de ventas y marketing, así como el equipo de
desarrollo y soporte.

Okticket es una de las empresas que más crecen tanto del sector como del propio Parque
Tecnológico. Debido a este rápido crecimiento, el equipo está en constante desarrollo y
cambio, tanto aquí en España como en el resto de sedes. Este crecimiento se refleja
en la recepción de un gran número de galardones y reconocimientos.
\footnote{\href{https://www.linkedin.com/posts/okticket_okticket-en-el-especial-startups-de-forbes-activity-7140622980618903552-UGWK}{Okticket en el especial startups 2023 de Forbes (LinkedIn)}}
\footnote{\href{https://www.elcomercio.es/economia/arcelor-okticket-premios-20230222002438-ntvo.html}{Arcelor y Okticket, premios nacional de Ingeniería Informática (EL COMERCIO)}}
\footnote{\href{https://www.okticket.es/blog/empresa-pyme-innovadora}{Okticket recibe el sello Pyme Innovadora (okticket.es)}}
\footnote{\href{https://www.okticket.es/blog/okticket-empresa-emergente-certificada}{Okticket, empresa emergente certificada (okticket.es)}}

La parte principal del negocio es el núcleo del software como servicio (Software as a
Service en inglés, en adelante \textit{SaaS}), es decir, la aplicación completa tanto
para administradores como para empleados. Este SaaS se oferta a empresas de cualquier
tamaño, cuyo precio final varía en función del número de usuarios, las características
e integraciones que requiera la empresa cliente y el soporte que se ofrezca.

Recientemente se han añadido nuevas propuestas a la cartera de servicios ofertada por
Okticket, como la OKTCard {-} una tarjeta inteligente que gestiona automáticamente los gastos,
así como la inclusión de nuevos ``módulos'' de gestión de gastos y viajes.

Debido a todo este crecimiento, la empresa maneja una gran cantidad de datos importantes que se
encuentran en diferentes formatos y en diferentes ubicaciones, lo que dificulta su análisis y
explotación. Por otra parte, se depende de la consulta manual o de servicios de terceros para
poder analizar estos datos, lo que supone un coste adicional.

\section{Objetivos}\label{sec:objetivos}
El objetivo del proyecto es el desarrollo y despliegue de una infraestructura de datos que permita
la integración, almacenamiento y análisis de grandes volúmenes de datos, provenientes de múltiples
fuentes y en diferentes formatos. La infraestructura de datos debe ser escalable, flexible y
robusta, para poder adaptarse a las necesidades cambiantes de la empresa.

La infraestructura de datos debe permitir la integración de datos de múltiples fuentes, tanto
internas como externas, y en diferentes formatos, como bases de datos, archivos de registros, APIs,
entre otros. La integración de datos debe ser automática y programable, para poder automatizar el
proceso de ingestión de datos y reducir el tiempo y los costes asociados.

Pese a que el entregable principal de este proyecto es la creación de una infraestructura, se
esperan también otros entregables en forma de herramientas de software, como \textit{scripts},
que faciliten la integración y análisis de los datos, así como la visualización de los mismos.

\chapter{Resultados y trabajo futuro}
El propósito de este capítulo es presentar las conclusiones obtenidas a partir
del desarrollo del proyecto, recopilar las dificultades encontradas y proponer
líneas de trabajo futuro en vista a la amplicación y mejora del sistema.

\section{Resultados}
El resultado del proyecto es un sistema de monitorización y análisis de datos
funcional y escalable, que permite a los usuarios obtener insights valiosos a
partir de la información recopilada.

Se ha logrado implementar una arquitectura robusta y flexible, basada en
tecnologías modernas y en la nube, que facilita la ingesta, procesamiento y
visualización de datos de manera eficiente y sencilla.

El sistema es capaz de ingestar datos de diversas fuentes, como bases de datos
internas, logs de aplicaciones y servicios externos, y de presentarlos de forma
clara y útil a través de Kibana.

La infraestructura se despliega y orquesta de manera automática en la nube,
permitiendo una gestión sencilla y eficiente del sistema para los
administradores.

Pese a que no se han completado todas las historias de usuario planificadas,
se han logrado los objetivos principales del proyecto, sentando las bases para
futuras iteraciones y logrando así entregar un producto mínimo viable (MVP).

El motivo de no haber completado todas las historias de usuario planificadas
se debe al tamaño y naturaleza del proyecto, que ha resultado ser más complejo
de lo esperado en un principio. Sin embargo, se ha logrado superar los objetivos
principales y se ha entregado un producto funcional y de calidad.


\newpage{}
\section{Trabajo futuro}
El proyecto ha sentado las bases para un sistema de monitorización y análisis de
datos robusto y escalable. Sin embargo, existen áreas de mejora y ampliación que
podrían ser abordadas en futuras iteraciones.


\subsection{Historias de usuario restantes}
Lo primero de todo sería completar las historias de usuario menos críticas que
han quedado pendientes, como la ingesta de datos de APIs externas o el
\textit{web scraping}. Estas funcionalidades permitirían enriquecer los datos
disponibles y ampliar las fuentes de información.

\begin{table}[H]
	\centering
	\begin{tabular}{|p{0.7\linewidth}|c|c|}
		\hline
		\textbf{Nombre} & \textbf{Prioridad} & \textbf{Tamaño} \\
		\hline
		\hline
		Como desarrollador de Okticket, quiero que los datos contengan metadatos que faciliten su filtrado o búsqueda & P2\cellcolor{yellow!50} & XS\cellcolor{blue!25} \\
		\hline
		Como trabajador de Okticket, quiero poder ver y consultar datos de empresas cliente & P2\cellcolor{yellow!50} & S\cellcolor{green!25} \\
		\hline
		Como gestor de una empresa cliente, quiero poder ver información relevante sobre mi empresa que recoja Okticket & P2\cellcolor{yellow!50} & M\cellcolor{yellow!50} \\
		\hline
		Como desarrollador de Okticket, quiero poder ingestar datos de APIs externas a la empresa & P2\cellcolor{yellow!50} & M\cellcolor{yellow!50} \\
		\hline
		Como desarrollador de Okticket, quiero poder ingestar información de páginas web externas (\textit{scraping}) & P2\cellcolor{yellow!50} & S\cellcolor{green!25} \\
		\hline
	\end{tabular}
	\caption{Historias de usuario restantes para futuras iteraciones}
	\label{tab:remaining_tasks}
\end{table}

Como se puede observar en la tabla \ref{tab:remaining_tasks}, las historias de
usuario restantes son las menos prioritarias (siguiendo con
\fullref{sec:planif_inicial}).

\newpage{}
\subsection{Integración de lenguaje natural para búsqueda (DSL)}
Sería interesante explorar la posibilidad de integrar el sistema con
un sistema de búsqueda y filtrado de datos a través de lenguaje natural, como
\textit{Natural Language Processing} (NLP)\footnote{
	\url{https://www.elastic.co/guide/en/elasticsearch/reference/current/query-dsl-query-string-query.html}
}

Esto permitiría a los usuarios realizar consultas de manera más intuitiva y
eficiente, sin necesidad de conocer la sintaxis de Kibana Query Language (KQL).


\subsection{Aplicación de modelos de Lenguaje de Gran Escala (LLM)}
Desde la empresa, se ha propuesto la posibilidad de aplicar modelos de LLM
para la generación de texto automática, presumiblemente de manera agnóstica al
proveedor (OpenAI, Anthropic, Google\ldots).

Esto permitiría la generación de informes y análisis de manera automática a
partir de los datos recopilados, facilitando la toma de decisiones y la
comunicación de insights a los usuarios.

\subsection{Más perspectivas futuras}
Gracias a la flexibilidad del stack ELK, se podrían añadir nuevas fuentes de
datos y visualizaciones, como logs de otras partes de la aplicación (gestor web,
otros servicios internos como Hubspot o Holded, aplicación móvil\ldots) o
visualizaciones más avanzadas y personalizadas.

Lo bueno de haber diseñado la arquitectura de la manera que se ha hecho
es que se pueden añadir nuevas funcionalidades sin necesidad de modificar
la infraestructura existente, simplemente añadiendo nuevos servicios y
configuraciones.

La escalabilidad del sistema también es un punto a tener en cuenta. En caso de
necesitar más capacidad de procesamiento o almacenamiento, se podría establecer
un sistema de autoescalabilidad en base a las definiciones ya realizadas.

Por último, se podría explotar la funcionalidad del stack de generar alertas
en base a la información ingestada y procesada, para notificar a los usuarios
de eventos críticos o anomalías detectadas en los datos.


\newpage{}
\section{Conclusiones y retrospectiva}
El desarrollo de este proyecto me ha permitido, a nivel personal, adquirir
conocimientos y habilidades muy útiles en áreas que tienen una gran demanda en
el mercado laboral actual, como la ingesta y procesamiento de datos, la
orquestación de servicios en la nube y la experiencia con proveedores cloud como
AWS.

A nivel técnico, he aprendido a trabajar con tecnologías modernas y potentes,
como Kafka, Logstash, Elasticsearch y Kibana, que me han permitido implementar
una solución robusta y escalable para la monitorización y análisis de datos.

Además, he tenido la oportunidad de trabajar con metodologías ágiles y de
gestión de proyectos, como Scrum, en proyectos reales con repercusiones y
resultados tangibles.

En cuanto a la retrospectiva del proyecto, considero que el tiempo dedicado ha
merecido la pena, pero podría haber sido gestionado u organizado diferente.


%% Esto incluirá la bibliografía correctamente en nuestro trabajo
\newpage % En una nueva página
\addcontentsline{toc}{chapter}{Bibliografía} % Añade la referencia al índice de contenido

\bibliographystyle{ieeetr} % Define el estilo de la bibliografía
\bibliography{biblio} % Indica el archivo que contiene la colección de citas

\nocite{template}

\end{document}
